Dataset-based evaluation cannot expose network-level failure modes. Intrusion detection systems for the Internet of Medical Things (IoMT) are trained and scored on stored captures, where a detector has no queue to slow down, no legitimate device to disconnect, and no feedback from its own decisions. We present a five-step methodology for placing a published IoMT intrusion detector inside a network simulator: obtain the system, audit its reported numbers arithmetically, rebuild it under leakage controls, place its decision rule in the data path of an ns-3 fog network, and evaluate it across 20 seeds, a scenario matrix, a sensitivity analysis and measured inference costs. Applying it to a recent fog-based system and two further IoMT detectors from 2025 and 2026 shows the following. Re-derivation from published confusion matrices confirms a positive-class inversion and transposed precision and recall in one paper and a train/test split drawn after oversampling in two of three; in one of them the test set holds more attack records than the whole dataset. On a public IoMT dataset, oversampling before the split alone raises F1 from 0.673 to 0.959 over 20 partitions, and duplicate rates depend on the feature space retained, from 0.52\,\% of raw rows to 98.95\,\% of modelling vectors on Edge-IIoTset. In the network, three threshold rules block every attacker in the baseline yet quarantine between 45\,\% and 100\,\% of legitimate wearables, and the published flat rule quarantines all of them in every one of 16 scenarios. A congestion-coupled threshold latches off under load in 20 of 20 baseline seeds and in 11 of 15 further scenarios at the same cost; it also fails at an otherwise safe cost when one fog node carries the flood, and on a Raspberry Pi 4 running the trained ensemble it latches at the Review~2 attack scale even at the board's idle inference cost (3.0\,ms per window). Three mitigations each restore full detection in every run, but on the Pi wearable telemetry still waits 0.3 to 1.2\,s behind each flood. A trade-off analysis shows that threshold adaptation slides along a curve set by the information it acts on, while gating it with a flow classifier of low false-positive rate moves the curve, cutting false alerts by a third with no missed attacker. We also report that a three-branch ensemble does not beat its best branch, that attack types withheld from training go almost undetected, and that the requested cross-dataset transfer is infeasible because the two datasets share no features. Code, raw runs and analysis are released.
